\subsection{Magnetic Catalysis and Fermion Mass Generation in de Sitter Spacetime --- arXiv:2608.07270}

\textbf{Authors:} Kohei Fujikura, Toshifumi Noumi, Shintetsu Yamazaki

\paragraph{主な主張}
de Sitter 時空中の一様磁場下で荷電フェルミオンのモード関数を Bunch--Davies 真空条件から構成し、NJL 模型の平均場近似でギャップ方程式を解くことで、磁場は磁気触媒によりカイラル対称性の破れを促進する一方、de Sitter 曲率は対称性回復を促し、両者の競合による相構造が生じることを示した \cite{Fujikura:2026magnetic}。

\paragraph{新規性}
小曲率展開に依存せず de Sitter 曲率と一様磁場を同時に扱い、強磁場・大曲率極限の解析式と数値的な相図を組み合わせて magnetic catalysis と curvature-induced symmetry restoration の競合を定量化した点が新しい。

\paragraph{位置付け}
インフレーション期の強磁場がフェルミオンの動的質量生成やカイラル相転移に与える影響を調べる基礎解析であり、NJL 模型を出発点としてより現実的なゲージ相互作用や inflationary magnetogenesis へ接続するための基準となる研究である。
